\section{Introduction}
\subsection{Problem and motivation}
A BallBot balances a body above a single spherical wheel and must deliberately alter its attitude to move. Transporting a liquid-containing payload introduces internal motion into this balancing task. An acceleration that reduces position error may excite the payload, while a reaction associated with payload motion changes the mechanical state on which the balance controller acts. Trajectory tracking, stabilization, and internal-energy management are therefore coupled objectives. A controller cannot be judged by liquid energy alone if it reduces that energy by remaining near the start of a commanded path.

The foundational BallBot work of Lauwers, Kantor, and Hollis established an inverse mouse-ball actuation architecture and experimentally demonstrated balancing and locomotion \cite{lauwers2006}. Those results motivate the robot class studied here; they do not establish the parameter fidelity of the smaller OpenBallBot-RL model. The present investigation uses a repository-specific multibody representation and records its geometry, mass distribution, actuator constraints, and contact assumptions explicitly. No physical BallBot was operated for the experiments analyzed in this manuscript.

The target task is a fixed figure-eight reference over mildly rough terrain. It combines longitudinal and lateral motion with changing reference curvature. The roughness is a procedural heightfield, rather than a collection of arbitrary obstacles or a measured outdoor surface. Calling the environment mildly rough is consequently important: results under these represented surfaces do not establish capability on steps, loose material, or discontinuous terrain. The liquid state is represented by a passive modal mechanical surrogate embedded in the same multibody plant. It is dynamic and produces reaction effects, but does not resolve a fluid surface.

This study addresses three research questions. RQ1 asks how BallBot locomotion, guarded terrain contact, and a reduced-order liquid surrogate can be integrated in a reproducible native simulation. RQ2 asks what the recorded classical and model-based evaluations establish about tracking, balance, and modal-energy trade-offs. RQ3 asks what privileged teachers and bounded PPO residuals can add to a strong model-based controller, and which observed limitations prevent a general advantage claim. Each question is answered at the scope of the archived experiments rather than by extrapolating a favorable development result.

\subsection{Why evidence boundaries matter}
The research archive contains experiments with different purposes. A short causal perturbation tests whether a particular action can improve a local outcome. A rolling teacher forecast evaluates a candidate plan from a saved state. A training episode supplies on-policy experience. A full tracking benchmark checks the original task contract. These records share a simulator, but their experimental units, time horizons, and selection mechanisms differ. Pooling them into one success-rate table would obscure what was actually tested.

A second difficulty is simulation validity. A numerically finite trajectory can still be invalid because contact geometry is unsupported, a liquid modal coordinate exceeds its declared approximation domain, or action logging associates a command with the wrong observation. Conversely, failure of a relative tracking-performance constraint does not necessarily mean that the mechanical system has become unstable. The archive includes both classes of issue. This paper preserves their distinction and retains unfavorable outcomes, including a known full-duration tracking failure and short teacher energy regressions.

A third difficulty is the meaning of privileged learning. True simulated liquid coordinates can help a modal-feedback controller, a planning teacher, or a residual actor. However, supplying only robot observations to one learned component does not make the complete controller deployable if its frozen backbone still consumes liquid state. We therefore describe information access at the whole-system level. Hardware transfer and causal liquid-state estimation are separate unresolved problems, rather than implicit consequences of successful policy inference.

\subsection{Scope and contributions}
The contribution is an evidence-based simulation study, not a new general theory of fluid dynamics or optimal control. First, the implementation integrates a source-audited modal payload into native robot dynamics, documents mass bookkeeping and energy conventions, and exposes the tested contact domain. This is an integration of established mechanical modeling and simulation tools, with repository-specific engineering and verification.

Second, the experimental methodology separates task completion, physical and numerical admissibility, and performance relative to a reference controller. It preserves historical protocols while constructing an experiment registry, comparability matrix, and claim-to-artifact mapping. This contributes a reproducible way to interpret heterogeneous iterative robotics experiments without recasting partial trials as completed tasks.

Third, the controller evidence establishes that the privileged modal-feedback configuration can complete recorded full-horizon cases while remaining sensitive to the initial terrain context. Short teacher evaluations demonstrate that the sign of a modal-energy effect can depend on the averaging window, even when trajectories share an identical prefix. These are empirical findings within a finite represented domain, not universal robustness or suppression guarantees.

Fourth, a matched exploration experiment tests an exact-likelihood motor-axis Gaussian mixture under unchanged controller authority and PPO settings. It distinguishes increased action-space coverage from improved closed-loop learning. The negative and mixed outcomes are part of the contribution: more probability near useful local actions did not yield consistent joint tracking and energy superiority over Standard PPO in the available replications.

The paper deliberately narrows claims where the archive is insufficient. It does not offer a completed, statistically balanced full-horizon comparison of PID, teacher, BC, and PPO. It also does not substitute the modal study for the project's longer-term resolved-fluid objective. Internal equation checks, native force accounting, and regression tests support implementation consistency; independent liquid and hardware measurements would be needed to assess physical predictive accuracy.

The remainder reviews relevant foundations, specifies the implemented plant and controllers, defines the frozen protocols and statistics, presents separately scoped results, and discusses limitations and reproducibility. Complete numerical provenance is provided in the accompanying evidence registries. This organization allows successful cases, censored experiments, and negative learning results to coexist without imposing a single unsupported ranking.
